% ============================================================
% WearLens report -- Sections 6-8 (revision 2, after AI review round 3)
% Figures: figs/interface.png, figs/scenario1_col.png, figs/scenario2_col.png
%   (scripts/report_figs.js + scripts/compose_report_figs.py).
% Every number is traceable to data/verify_expected.json and was checked in the browser (scripts/verify_app.js).
% ============================================================
\providecommand{\wlfig}[2]{%
  \ifdefined\includegraphics
    \IfFileExists{#1}{\includegraphics[width=#2]{#1}}{\fbox{\parbox{0.9#2}{\centering Missing file: #1}}}%
  \else\fbox{\parbox{0.9#2}{\centering Load graphicx to show #1}}\fi}

\section{System Design}
\begin{figure*}[t]
\centering
\wlfig{figs/interface.png}{\textwidth}
\caption{The WearLens interface. Top: guided steps and global controls (feature, normalization, VB limit,
layer-flag thresholds, machine legend). V1~wear curves, V2~replacement rules, V3~correlation ranking,
V4~feature against wear, V5~run timeline of one end mill, V6~details of one run. Enlarged panels of V2, V5 and
V6 are shown in Figs.~\ref{fig:s1} and \ref{fig:s2}.}
\label{fig:interface}
\end{figure*}

WearLens is a single page with six coordinated views and one row of global controls (Fig.~\ref{fig:interface}).
Three guided steps set the controls for CH1--CH3 and show a guiding question; every setting remains editable.
In V1, V2 and V4--V6 \emph{colour} encodes only the machine, with three hues whose worst pairwise distance under
simulated colour-vision deficiency is $\Delta E_{\mathrm{OK}}=9.2$ (OKLab $\times100$; 8 is a common target).
Trend lines add a line style per machine, and machines are always labelled in text. End mills are not coloured,
because nine hues cannot be told apart; they are identified by position and labels. V3 uses its own diverging
scale for $\rho$ and shows no machine colours. Machine-control features are listed first; dynamometer features
are greyed and marked ``reference only''. Quantities with different units are stacked in separate panels
instead of dual axes. The views use different horizontal axes on purpose (contact time in V1 and V2, VB in V4,
run in V5); they are linked by highlighting the same runs and end mills rather than by a shared axis.

\paragraph{V1 Wear curves (T1).}
Small multiples ($3\times3$, one row per machine, shared scales) show VB against contact time. Position on a
common scale compares trends most accurately; aligned rows let the eye compare the wear phases between
machines, while each panel stays selectable and brushable for one end mill. An overlay
(Fig.~\ref{fig:motivation}a) compares machines more directly but makes nine lines overlap and hides
individual end mills, so we use it only in the report.

\paragraph{V2 Replacement rules (T2, T6, T7).}
Each end mill is a horizontal bar from zero to the contact time at which VB reaches the limit, grouped by
machine. A bar encodes this interval and splits it into the part used before replacement (saturated) and the
unused part (light) or overrun (hatched); a dot plot would show only the end points. Fixed times are set by a
draggable line or, in leave-one-out mode, by one tick per end mill. In the signal-alarm mode, a triangle marks
the alarm, the light bar is the remaining life until the limit, and a label states the error; a summary reports
MAE and the numbers of early, late and missed alarms.

\paragraph{V3 Correlation ranking (T3).}
A heatmap shows Spearman $\rho$ between each feature (rows) and VB per end mill and pooled (columns),
with a diverging blue--grey--red scale centred at $\rho=0$ and printed values. A matrix reveals
machine-specific blocks that a bar chart of one column would hide. Rows are sorted by pooled $\rho$ and move
with an animated transition when the ranking changes, so the user can follow a feature.

\paragraph{V4 Feature against wear (T4, T7).}
A scatterplot shows the normalized feature of all 6,418 runs against VB at low opacity, overlaid with a
rolling median per end mill. A band shows the range of the per-end-mill level at the VB limit, from which
alarm levels are calibrated. Hexagonal binning would hide machine identity.

\paragraph{V5 Run timeline (T5, T6).}
Dots show the feature per run of one end mill, alternating bands the layers, and badges the flagged layers
(J, R); the header counts flagged layers over all end mills. In signal-alarm mode, a black line shows the
trailing median an online monitor would see, a dashed line the alarm level, and a triangle the alarm.
A timeline is the only view in which layer structure is visible; plotting against VB (as in V4) would compress
the many runs that share an interpolated VB value.

\paragraph{V6 Run details (T6).}
Stacked panels with a shared time axis show 200-bin envelopes of spindle torque, feed-drive effort and the
dynamometer force (reference) for one run; a second run can be pinned in grey. Separate panels avoid a dual
axis and let the user check whether independent signals change together; a 20\,ms raw force snippet shows
the individual cutting edges.

\paragraph{Coordination.}
Brushing a time range in V1 recomputes V3 and fades other runs in V4 and V5; brushing in V4 highlights the
runs in V1 and V5; clicking a row of V3 changes the feature in V2, V4 and V5; selecting an end mill anywhere
highlights it everywhere; clicking a run in V4 or V5 opens it in V6. Changing the normalization, VB limit or
calibration updates V2--V5.

\begin{table}[t]
\centering\small
\caption{Task--design mapping. Every view and control serves at least one task.}
\label{tab:mapping}
\begin{tabular}{@{}p{0.09\linewidth}p{0.34\linewidth}p{0.45\linewidth}@{}}
\toprule
Task & Views & Interactions / analysis \\
\midrule
T1 & V1 & select end mill, VB limit, step 1 \\
T2 & V2 (fixed time), V1 & drag time, per-machine and leave-one-out modes, VB limit; A3 \\
T3 & V3 & normalization, V1 time brush, step 2; A1, A2 \\
T4 & V4, V3 & select feature, normalization, brush; A1 \\
T5 & V5 & jump and rise thresholds, step 3; A4 \\
T6 & V5, V6, V2 & select and pin runs, dynamometer toggle; A3, A4 \\
T7 & V2 (signal alarm), V4, V5 & calibration, alarm level, persistence; A1, A3 \\
\bottomrule
\end{tabular}
\end{table}

\section{Analysis Components}
All components run in the browser, so every parameter can be changed and all views update immediately.
\begin{description}
\item[A1 Normalization.] The three normalizations of Section~\ref{sec:users} are offered side by side and
labelled ``production'' or ``trial only'', so the user sees whether a feature's quality depends on information
that will not exist in production (T3, T4, T7).
\item[A2 Correlation ranking.] Spearman $\rho$ is recomputed per end mill and pooled whenever the
normalization or time range changes (T3). It screens features; it does not validate them.
\item[A3 Replacement-rule evaluation.] Fixed times are one per fleet or per machine (default: the earliest
time at which an end mill of the group reaches the limit; draggable), or leave-one-out (each end mill gets the
earliest life of the \emph{other} end mills of its machine). The signal alarm compares the trailing median of
15 runs with an alarm level that is set manually or calibrated on the other machines, on the other end mills of
the same machine (leave-one-out) or in-sample; it fires after the chosen persistence. Both rule types report
the error of every end mill, MAE, and the numbers of early, late (overrun) and missed end mills, so that rules can
be compared under the same test (T2, T6, T7).
\item[A4 Layer diagnostics.] The \emph{jump} of a layer is the ratio of its median to the mean of the two
neighbouring layer medians minus one; positive and negative jumps are flagged by their absolute value, and
only the local maximum is flagged because a jump also appears in the neighbours. The \emph{rise} is the change
of a linear fit across the layer relative to the layer median. Both use the raw feature and user-set
thresholds (T5, T6).
\end{description}
Together, A1--A4 let the user state a hypothesis (``a signal rule replaces fixed times'', ``this early alarm is
caused by one layer'') and test it in the same session.

\section{Usage Scenarios and Insights}
The engineer uses the spindle torque, the production normalization and a VB limit of 140\,$\mu$m.

\paragraph{Scenario 1: Fixed time or signal alarm? (CH1, CH2).}
In V1 the M2 end mills differ from the others mainly in the slow middle phase (T1). In V2, one fixed time for
all machines (53.5\,min, set by the shortest-lived end mill) leaves 34--36\,\% of the life of each M2 end mill
unused, and one time per machine reduces the mean unused life to 1.9\,\% (T2). Because both times are fitted
to the end mills they are evaluated on, the engineer switches to the leave-one-out mode
(Fig.~\ref{fig:s1}a): the MAE is 1.7\,min, with three overruns of at most 2.0\,min and no missed end mill.
In V3 the spindle torque has $\rho\ge0.97$ in every end mill and 0.95 pooled, against 0.57 with raw values
(T3), and V4 shows the three machines rising in the same direction (T4). In V2 a level calibrated on the other
two machines alarms 5.8--14.7\,min early on M1 and M2 and misses all three M3 end mills, so alarm levels do
not transfer between machines. Calibrated leave-one-out on the same machine (Fig.~\ref{fig:s1}b), the alarm
reaches an MAE of 4.7\,min over eight alarms, with one late alarm (E3, 2.7\,min) and one missed end mill (E5)
(T7). \textbf{Insight:} under the same leave-one-out test, a fixed time per machine beats the spindle-torque
alarm, because the lives of the end mills on one machine differ by less than the alarm error. For this tool and
process the trial supports per-machine fixed times; a signal alarm would need machine-specific calibration
and would serve at most as a safety net. Both conclusions are limited by three end mills per machine and by
VB measurements every 2.4--4.7\,min.

\begin{figure}[t]
\centering
\wlfig{figs/scenario1_col.png}{\linewidth}
\caption{Scenario 1 in V2, both rules tested leave-one-out. (a) Fixed time per machine: dashed ticks are
the replacement times. (b) Spindle-torque alarm: triangles are alarms; light bars are the remaining life until
VB\,=\,140\,$\mu$m.}
\label{fig:s1}
\end{figure}

\paragraph{Scenario 2: How one event shapes the rule (CH3, CH1).}
The largest alarm error in Scenario~1 belongs to E1 (12.3\,min early at VB\,=\,103\,$\mu$m), under every
calibration. V5 shows that the alarm is raised in layer 15, one of only two flagged layers among 222 at the
default thresholds (the other is the first layer of E3, during break-in) (T5; Fig.~\ref{fig:s2}). The engineer
pins run 420 of layer 14 and opens run 443 in V6 (T6): spindle torque and dynamometer force are higher over
the whole run, while the feed-drive effort barely changes. Two independent chains agree, so a sensor fault is
unlikely; the torque falls back in layer 16, so it is not wear. To absorb such an event, the persistence must
exceed a whole layer: with the level calibrated on the other machines and 29 runs, the alarm is still
9.9\,min early, because all runs of the layer are high; with 30 runs it moves to 0.5\,min early. The same
30-run persistence, however, leaves five of nine end mills without an alarm under in-sample calibration. The
engineer therefore keeps a persistence of one run and adds a rule: an alarm raised in a flagged layer
triggers an inspection, not a replacement. The early alarms of E2--E9 do not fall into flagged layers, so they
are caused by the alarm level, not by events; without E1, the leave-one-out MAE is 3.6\,min, still above the
1.7\,min of the fixed time. \textbf{Insight:} a single layer-boundary event dominated the evaluation of the
signal rule; identifying it showed that the rule's remaining weakness is the alarm level itself, which confirms
the conclusion of Scenario~1.

\begin{figure}[t]
\centering
\wlfig{figs/scenario2_col.png}{\linewidth}
\caption{Scenario 2. V5: E1 with trailing median (black), alarm level (dashed) and alarm (triangle) in layer 15,
flagged J. V6: run 443 (layer 15) against the pinned run 420 (layer 14, grey).}
\label{fig:s2}
\end{figure}
